\section{Amortising the RL policy into the model}
\label{sec:loop}
A pathwise-trained policy improves most under the sampler it was trained with and less under cheap samplers (the 64-step ODE, $\gamma4$@16). The gap is closed with the model's own verified samples. One \emph{round} has three steps: sample the current policy under the deployment sampler on a fixed pool of training-side problems, several seeds per problem; keep the samples the verifier accepts, at most four per problem (one on the unique-solution tasks); fine-tune the model on them with the base's supervised objective, starting from the current policy. Rounds are chained, each from the previous round's final checkpoint, for a fixed number of rounds per task, and the last round is reported. The 2D critic gives the loop a second use: because its field scores partial states, the pool can be sampled with candidates pruned at an intermediate step (best of four) and pushed along the field's gradient, which restores the pool's coverage when plain sampling stalls; Countdown's last two rounds use this critic-guided sampler (Appendix~\ref{app:cd_loop}). The loop is rejection-sampling fine-tuning \citep{yuan2023rft,singh2023restem,zelikman2022star} on a continuous-state model, and any policy can enter it; \S\ref{sec:experiments} asks where its gains come from with three controls that hold the problems and the rounds fixed: the identical loop from the SFT base, the identical loop from the GRPO policy, and supervised fine-tuning on the \emph{gold} derivations of exactly the problems the loop sampled.
